\chapter*{Conclusion Générale et Perspectives}
\addcontentsline{toc}{chapter}{Conclusion Générale et Perspectives}

Les agents d'intelligence artificielle autonomes représentent une évolution importante des applications basées sur les LLM. Leur capacité à utiliser des outils, à accéder à des données et à enchaîner des actions transforme une erreur de raisonnement ou une instruction malveillante en risque opérationnel. Dans ce contexte, la question de sécurité ne peut plus être formulée uniquement comme « le modèle produit-il une réponse sûre ? ». Elle devient : « l'application doit-elle exécuter l'action que le modèle vient de proposer ? ».

Le projet \Aegis{} a été conçu autour de cette distinction. Son principe central, \textbf{« le LLM propose, Aegis décide »}, sépare le moteur de raisonnement de l'autorité de sécurité. La cible architecturale introduit un point de contrôle obligatoire, \Guard{}, capable d'évaluer le but utilisateur, l'identité, la provenance, le niveau de confiance du serveur MCP, la classification de la donnée, le risque de l'outil et la destination. Cette approche vise à fournir une défense qui ne dépend pas uniquement de la capacité du modèle à reconnaître une attaque textuelle.

Le projet ne se limite pas à la défense. \Red{} apporte une méthodologie de test reproductible destinée à générer et rejouer des scénarios d'injection, d'empoisonnement, d'abus de capacités ou de supply chain MCP. \GRC{} complète la chaîne en transformant les résultats techniques en risques, contrôles, preuves, indicateurs et décisions de traitement. Ce lien entre attaque, contrôle et risque résiduel constitue l'un des axes les plus importants du projet : une défense n'est pas seulement « présente », elle doit être testée et produire une preuve observable.

La réalisation effectuée à la date de cette version du rapport a permis de valider plusieurs briques fondamentales. Le provider OpenRouter fonctionne et a été durci pour gérer certaines défaillances temporaires. Le prototype reçoit des appels d'outils structurés et expose clairement la frontière entre proposition du LLM et exécution réelle. Une baseline volontairement non protégée, munie d'un outil de lecture et d'une messagerie simulée, a été construite. Un premier corpus Aegis Red de prompt injections indirectes a été exécuté avec instrumentation des accès et de l'outbox.

Les résultats préliminaires ont une valeur méthodologique particulière. Les cinq variantes d'injection indirecte testées contre Nemotron 3 Ultra n'ont pas provoqué d'accès au fichier privé ni d'exfiltration. Il serait erroné d'en conclure que la menace n'existe pas : la littérature montre que le succès des attaques varie fortement selon le modèle et le scénario. Le résultat correct est que ce modèle, dans cette configuration, a résisté aux variantes testées. Cette observation incite à ne pas « fabriquer » une baseline vulnérable par des instructions artificiellement permissives, mais à construire une évaluation multi-modèles et multi-attaques.

Un deuxième enseignement concerne la dépendance aux services externes. Une surcharge NVIDIA a produit une erreur 502 récupérée par la logique de retry. Deux autres modèles gratuits référencés dans une première comparaison n'étaient plus disponibles et ont retourné des erreurs 404. Ces incidents montrent que la reproductibilité d'une expérimentation agentique exige d'enregistrer non seulement le code et les prompts, mais aussi le provider, le modèle, l'état de l'endpoint et les erreurs d'infrastructure. L'ajout ultérieur d'un modèle local figé constituera un complément pertinent.

Les perspectives immédiates du projet sont donc clairement identifiées. La première est l'intégration du SDK MCP Python v2 et la migration des fonctions locales vers des serveurs Documents et Messaging. La deuxième est l'implémentation d'Aegis Guard v1 avec des règles simples mais déterministes : par exemple, bloquer une donnée RESTRICTED vers une destination externe lorsque le but utilisateur ne justifie pas l'action. La troisième est l'ajout de provenance, de classification et d'identité. La quatrième est la construction d'un pipeline de télémétrie structuré, puis du modèle GRC.

Après cette phase, l'évaluation pourra être étendue à des scénarios spécifiquement MCP : serveur non approuvé, tool poisoning, tool shadowing, résultats d'outils malveillants et abus de scopes. Le RAG et la mémoire persistante permettront ensuite d'étudier les attaques qui survivent à travers le contexte. Enfin, un runner multi-modèles mesurera l'Attack Success Rate, le taux de prévention, les faux positifs, la réussite des tâches légitimes, la latence et le blast radius.

Sur le plan GRC, l'objectif final n'est pas d'afficher un pourcentage arbitraire de conformité. Il est de construire une chaîne de preuve transparente. Un risque doit pouvoir être relié à une attaque reproductible ; un contrôle à une implémentation ; l'implémentation à un test ; le test à une preuve ; la preuve à une évaluation d'efficacité ; et cette efficacité au risque résiduel. Les référentiels NIST AI RMF, ISO/IEC 42001, ISO/IEC 23894, certains contrôles ISO/IEC 27001 et le cadre marocain de protection des données fournissent des points de correspondance, mais AegisMCP ne prétend pas remplacer un audit de conformité officiel.

À terme, AegisMCP peut évoluer au-delà du cadre académique. La couche Guard peut devenir un \textit{policy enforcement point} réutilisable pour plusieurs agents. Le catalogue Red Team peut servir de test de non-régression avant le déploiement d'un nouveau modèle ou serveur MCP. Le tableau de bord GRC peut fournir aux équipes sécurité et gouvernance une vue commune sur les risques agentiques. Une extension pourrait également intégrer des politiques signées, l'attestation des serveurs MCP, des scopes de capacités temporaires, des approbations multi-niveaux et des politiques de flux de données plus fines.

Le projet défend finalement une idée simple mais structurante : \textbf{l'autonomie d'un agent IA ne doit pas être confondue avec une autorité illimitée}. Plus l'agent est capable d'agir, plus ses actions doivent être contextualisées, minimisées, contrôlées et auditées. AegisMCP cherche à transformer ce principe en architecture expérimentale mesurable, afin de contribuer à une utilisation plus sûre et gouvernable des agents IA connectés aux outils du monde réel.
